% Part II --- Bargaining (draft Aulas 2-3). Numbers checked by check_part2.py.

\textbf{Read.} \MR{} Ch.~2, ``Wealth Effects, Value Maximization and the Coase Theorem''
(pp.~35--39); Ch.~5, ``Achieving Commitment'' (pp.~133, 139--140), ``Private Information and
Precontractual Opportunism'' and ``Measurement Costs and Investments in Bargaining Position''
(pp.~140--149). Varian (2014, with Calculus) \S30.7 ``Bargaining'' and ``The Ultimatum Game''
(pp.~581--583), \S35.2 (p.~670). MWG Ch.~9, Appendix~A (pp.~296--299). Cabral (2017) \S5.1
(pp.~96--99).

In the draft schedule, Bargaining takes two classes: Aula~2 covers
\S\ref{sec:b-question}--\S\ref{sec:b-nash} (value maximization, Coase, the division of the
gains), and Aula~3 covers \S\ref{sec:b-rub}--\S\ref{sec:b-costs} (alternating offers,
commitment, private information, bargaining costs).

% ---------------------------------------------------------------------------------------
\section{The question: is bargaining indeterminate?}\label{sec:b-question}

Every organization runs on agreements: a supplier contract, a division of profits among
partners, a transfer price between divisions, a wage. Each agreement is reached by some form of
bargaining. \MR{} open the topic with the traditional view that ``the outcomes of bargaining among
small numbers of people are essentially indeterminate'': bargaining strength, credibility,
guile, patience and strategic insight combine with initial conditions to decide the outcome,
and predicting it ``was too much for economic analysis'' \pg{140}.

This Part gives \MR's answer in three steps.

\begin{enumerate}
  \item \textbf{What is created is determinate.} When preferences show no wealth effects, every
  efficient agreement chooses the actions that maximize the parties' total value (the
  \emph{value maximization principle}). Bargaining power decides only the money transfers (the
  \emph{Coase theorem}). Sections~\ref{sec:b-wealth}--\ref{sec:b-coase}.
  \item \textbf{Who gets what can be modelled.} The Nash solution and the Rubinstein
  alternating-offers game turn ``bargaining power'' into threat points, patience and the right
  to propose. Under complete information these models predict immediate, efficient agreement.
  Sections~\ref{sec:b-nash}--\ref{sec:b-commit}.
  \item \textbf{Private information and bargaining costs break efficiency.} When each side knows
  something the other does not, the full-information efficient outcome can be infeasible, and
  even the best feasible arrangement gives up some gains. Time, haggling and information spent
  on shifting the division are wasted. Sections~\ref{sec:b-private}--\ref{sec:b-costs}.
\end{enumerate}

Throughout, a decision has two parts: $y$, the \emph{real} choice (what is produced, who owns
an asset, whether a project is built), and $x = (x_1,\dots,x_N)$, the money transfers among the
parties.

% ---------------------------------------------------------------------------------------
\section{Wealth effects and the value index}\label{sec:b-wealth}

A \emph{wealth effect} is a change in choices caused by a change in wealth: a poorer person may
be unable to afford an option, or may trade off risk differently \pg{35}.

\begin{defn}{No wealth effects \pg{35}}
A decision maker shows \emph{no wealth effects} with respect to a set of decisions $Y$ when, for
any two decisions $y_1, y_2 \in Y$:
\begin{enumerate}[label=(\roman*)]
  \item there is a definite amount of money $\Delta x$ that would just compensate him or her for
  a switch from $y_1$ to $y_2$;
  \item that amount does not change if he or she is first given extra wealth;
  \item he or she has enough money to absorb any wealth reduction needed to pay for a switch
  from the less preferred to the more preferred option.
\end{enumerate}
\end{defn}

\begin{prop}{Proposition 1 (the value index, \MR{} p.~36 and n.~7)}
Under no wealth effects, preferences over (money $x$, decision $y$) can be represented by
\[
  u(x, y) = x + v(y),
\]
where $v(y)$ is the cash-equivalent value of $y$. The compensating transfer for a switch from
$y_1$ to $y_2$ is $\Delta x = v(y_1) - v(y_2)$, the same at every initial wealth $M$.
\end{prop}

\emph{Why.} With $u = x + v(y)$, the person is indifferent between $(M, y_1)$ and
$(M + \Delta x, y_2)$ exactly when $M + v(y_1) = M + \Delta x + v(y_2)$, so
$\Delta x = v(y_1) - v(y_2)$: it exists (condition i) and does not depend on $M$ (condition ii).
Condition (iii) holds as long as $M$ exceeds the largest difference $|v(y) - v(y')|$. Conversely,
\MR's Q~2.4 (p.~53) asks you to show that conditions like (i)--(ii) force this form.

\textbf{Where it fails.} \MR{} contrast two cases \pg{35--36}. A corporation that finds itself
richer after a rise in its asset values would plausibly keep the same prices and hurdle rates:
no wealth effects. A worker who wins the lottery buys a house and quits: large wealth effects.
Condition (iii) fails when the transfer is large relative to the person's resources, as with a
nuclear-plant worker made financially liable for his or her mistakes \pg{36}. The rule of thumb:
the assumption is least safe ``when the decision makers are individuals and when large cash
transfers or significant changes in personal living conditions are involved'' \pg{36}.

\begin{example}{With wealth effects, who holds the right changes what is done}
A plant gains $50$ (in money) from dumping waste in a river. A resident with wealth $w = 100$ has
utility $u = \ln x + a\,y$, where $y = 1$ if the river is clean and $a = 0.5$. This is not of
the form $x + v(y)$: the money value of a clean river depends on wealth.

\emph{Willingness to pay} for a clean river (resident owns nothing, must buy it):
$\ln(w - P) + a = \ln w \;\Rightarrow\; P = w(1 - e^{-a}) \approx 39.35$.

\emph{Willingness to accept} for giving the river up (resident owns the right):
$\ln(w + A) = \ln w + a \;\Rightarrow\; A = w(e^{a} - 1) \approx 64.87$.

If the plant holds the right, the resident can offer at most $39.35 < 50$: the plant dumps. If
the resident holds it, the plant would have to pay at least $64.87 > 50$: the river stays
clean. The \emph{real} outcome $y$ depends on the assignment of rights. With wealth effects, the
gap between willingness to pay and to accept makes ``the efficient choice'' depend on who starts
with what, and the clean separation of the next section disappears.
\end{example}

% ---------------------------------------------------------------------------------------
\section{The value maximization principle}\label{sec:b-vmp}

\begin{prop}{The Value Maximization Principle (\MR{} p.~36)}
``An allocation among a group of people whose preferences display no wealth effects is efficient
only if it maximizes the total value of the affected parties. Moreover, for any inefficient
allocation, there exists another (total value maximizing) allocation that all of the parties
strictly prefer.''
\end{prop}

\textbf{Setting.} $N$ people with $u_n(x_n, y) = x_n + v_n(y)$. Decision $y$ generates net cash
$P(y)$ (positive for an investment, negative for a public good that costs money), and the
transfers must add up: $x_1 + \dots + x_N = P(y)$. The \emph{total value} of an allocation is
\[
  \sum_{n} \big[x_n + v_n(y)\big] \;=\; P(y) + \sum_n v_n(y) \;\equiv\; W(y),
\]
which depends on $y$ alone: changing the $x$'s moves value between people one for one and never
changes the total.

\begin{prop}{Proposition 2 (efficiency $\Leftrightarrow$ value maximization)}
An allocation $(y, x_1, \dots, x_N)$ is Pareto efficient if and only if $y \in \argmax_{y'} W(y')$.
\end{prop}

\begin{proof}
($\Leftarrow$) Let $y$ maximize $W$ and suppose another allocation $(y', x')$ made everyone at
least as well off and someone strictly better off. Summing utilities,
$W(y') = \sum_n [x'_n + v_n(y')] > \sum_n [x_n + v_n(y)] = W(y)$, contradicting the maximality of
$W(y)$.

($\Rightarrow$) Suppose $y$ does not maximize $W$: some $y^*$ has $G \equiv W(y^*) - W(y) > 0$.
Give each person the transfer
\[
  x^*_n = x_n + v_n(y) - v_n(y^*) + \frac{G}{N}.
\]
Then $u_n(x^*_n, y^*) = x_n + v_n(y) + G/N > u_n(x_n, y)$ for every $n$. The transfers add up:
using $\sum_n x_n = P(y)$ and $G = P(y^*) + \sum_n v_n(y^*) - P(y) - \sum_n v_n(y)$,
\[
  \sum_n x^*_n = P(y) + \sum_n v_n(y) - \sum_n v_n(y^*) + G = P(y^*).
\]
So every inefficient allocation is strictly Pareto dominated by a value-maximizing one, which is
the second sentence of the principle.
\end{proof}

\MR's Q~2.3 (p.~53) asks for this proof; the hint there is the equal-split construction used
above. The proof uses condition (iii): the transfers $x^*_n$ must be affordable.

\begin{figure}[ht]
\centering
\begin{tikzpicture}[scale=0.9,>=Stealth]
  \draw[->] (0,0) -- (6.2,0) node[below] {value to 1};
  \draw[->] (0,0) -- (0,6.2) node[left] {value to 2};
  \foreach \c/\lab in {2.5/{$W(y_1)$}, 4/{$W(y_2)$}, 5.5/{$W(y^*)$}} {
    \draw[thick, defc!70] (0,\c) -- (\c,0);
    \node[defc, font=\footnotesize, right] at (\c,0.25) {\lab};
  }
  \fill (1.6,2.4) circle (2pt) node[left] {$A$};
  \fill (2.15,3.35) circle (2pt) node[above right] {$B$};
  \draw[->, gray] (1.65,2.45) -- (2.1,3.28);
\end{tikzpicture}
\caption{\MR's Figure 2.1 (p.~37). Each line is the set of value pairs reachable with one
decision $y$ as the transfers vary; lines are at $45^\circ$ because transfers move value one for
one. Any point $A$ on a lower line is dominated by some $B$ on the highest line; points on the
highest line are undominated.}
\end{figure}

\begin{keyidea}
When the principle applies, ``the issue of the distribution of value is completely separable
from the issue of how value is created'' \pg{38}. Efficiency picks $y$; everything else
(bargaining power, patience, initial ownership) only moves the $x$'s.
\end{keyidea}

\textbf{The limit (\MR{} n.~10, p.~37).} With wealth effects, an allocation that maximizes the
sum of utilities is still efficient, but there can be efficient allocations that do not: value
maximization becomes \emph{sufficient} but not \emph{necessary} for efficiency. The utility
possibility frontier is then curved rather than a $45^\circ$ line, and different points on it
can use different $y$'s (as in the river example above).

\textbf{Range of application} \pg{38}. The variable $y$ can be the money spent on parks or
libraries, the identity of the person assigned overtime, the owner of water rights, the brand of
computer adopted, or the site of a new office building. In each case the principle says the
group should, and if it bargains efficiently will, pick the $y$ with the largest total value,
then settle shares by transfers.

% ---------------------------------------------------------------------------------------
\section{The Coase theorem}\label{sec:b-coase}

\begin{prop}{The Coase Theorem (\MR{} p.~38)}
``If the parties bargain to an efficient agreement (for themselves) and if their preferences
display no wealth effects, then the value-creating activities ($y$) that they will agree upon do
not depend on the bargaining power of the parties or on what assets each owned when the
bargaining began. Rather, efficiency alone determines the activity choice. The other factors can
affect only decisions about how the costs and benefits are to be shared ($x$).''
\end{prop}

The theorem is the value maximization principle plus one behavioural assumption: the parties
reach an agreement with no further mutual gains. Its organizational corollary is \MR's
foundation for the transaction-costs approach: for a given production plan, ``the efficient
organization for that plan is the one that minimizes transaction costs'' \pg{38}, where
transaction costs include writing contracts, supervising, enforcing and resolving disputes.

\begin{example}{The steel plant and the river (after Cabral, pp.~97--98)}
Cabral's setting: a steel plant can dump waste in a river, harming residents downstream; society
can give the right to clean water to the residents or the right to dump to the plant. Put
numbers on it: dumping is worth $100$ to the plant and costs the residents $60$, everyone has
no wealth effects, and bargaining is costless.

\begin{itemize}
  \item \emph{Plant holds the right.} The residents would pay at most $60$ to stop dumping; the
  plant needs at least $100$. No deal: the plant dumps, no money changes hands.
  \item \emph{Residents hold the right.} The plant can pay anything in $[60, 100]$ for
  permission. With an even split of the $40$ gain (the Nash solution of \S\ref{sec:b-nash}),
  the price is $60 + 20 = 80$. The plant dumps.
\end{itemize}
Dumping happens either way, because it creates $100 - 60 = 40$ of value. The assignment of rights
only decides whether residents end up $80 - 60 = +20$ or $0$, and the plant $+20$ or $+100$.
Reverse the numbers (harm $100$, gain $60$) and the river stays clean under either assignment.
\end{example}

Cabral draws the lesson that the Coase theorem is about efficiency, not fairness: ``From an
efficiency point of view, it makes no difference who owns the rights to a clean river (factory or
residents). This is not to say, however, that both property rights allocations are equally
fair'' (Cabral, p.~99). Varian's \S35.2 makes the same point in consumer theory: with
quasilinear preferences the efficient amount of the externality is independent of the
assignment of property rights, and only the distribution changes (Varian, p.~670).

\textbf{When the theorem fails.} Each hypothesis can break:
\begin{enumerate}
  \item \emph{Wealth effects and limited capital.} The implications ``do not hold when some of the
  parties have very limited capital with which to make payments''. \MR{} would apply the analysis
  to a contract between General Motors and Toyota, but not ``uncritically'' to land tenure in a
  developing country or to slavery in the pre-Civil War American South \pg{39}.
  \item \emph{Costly bargaining.} Haggling, delay and measurement absorb part of the gains
  (\S\ref{sec:b-costs}).
  \item \emph{Private information.} Efficient bargaining is an assumption of the theorem, not a
  result: with private values, no feasible procedure may reach the efficient $y$
  (\S\ref{sec:b-private}).
  \item \emph{Many parties.} Free riding makes unanimous agreement harder as the group grows
  (\S\ref{sec:b-many}).
\end{enumerate}

\textbf{Rival explanations of organization} \pg{39}. The transaction-costs view says
arrangements exist to raise the total wealth to be shared. Against the \emph{Marxian} view
(organization reflects power and class interests), \MR{} cite Yugoslavia, where worker-controlled
firms hired managers who organized factories much like capitalist firms. Against the
\emph{Harvard school} of industrial organization (arrangements raise market power), they read
resale price maintenance as a defence against discounters free riding on the service and customer
education provided by other dealers. They add that only systematic evidence can decide between
the views.

\begin{caution}{Three misreadings of Coase}
(1) ``Property rights do not matter.'' They matter for the distribution, and with wealth effects
for the allocation too. (2) ``Externalities are never a problem.'' The theorem assumes efficient,
costless bargaining; its practical lesson is to look for the transaction costs that block it.
(3) ``Coase predicts a fair outcome.'' It predicts an efficient $y$ and says nothing about $x$.
\end{caution}

% ---------------------------------------------------------------------------------------
\section{Dividing the gains: the Nash bargaining solution}\label{sec:b-nash}

The value maximization principle leaves the transfers open. Two families of models close them.
The \emph{axiomatic} approach asks which division satisfies reasonable properties; the
\emph{strategic} approach (next section) writes down the bargaining procedure as a game. Varian
introduces both and notes that the outcome of the Nash model depends ``on how risk averse the
players are and what will happen if no bargain is made'', leaving its full treatment out of his
book (Varian, p.~581). The statement below is the standard one (Nash 1950), not taken from the
course books.

\begin{defn}{Bargaining problem and Nash solution}
Two players can agree on a pair of payoffs $(u_1, u_2)$ in a feasible set $U$, or fail to agree
and get the \emph{threat point} (disagreement point) $d = (d_1, d_2)$. The Nash bargaining
solution is the unique rule satisfying Pareto efficiency, symmetry, invariance to affine
rescaling of utilities and independence of irrelevant alternatives; it selects
\[
  (u_1^N, u_2^N) \in \argmax_{(u_1, u_2) \in U,\; u_i \ge d_i} \; (u_1 - d_1)(u_2 - d_2).
\]
\end{defn}

With no wealth effects the frontier is the $45^\circ$ line $u_1 + u_2 = V$, where
$V = \max_y W(y)$ is the maximal total value. Maximizing $(u_1 - d_1)(V - u_1 - d_2)$ gives
\[
  u_1^N = d_1 + \tfrac{1}{2}\,(V - d_1 - d_2), \qquad u_2^N = d_2 + \tfrac{1}{2}\,(V - d_1 - d_2):
\]
each party gets its threat point plus half the \emph{surplus} $V - d_1 - d_2$. The
\emph{generalized} Nash solution maximizes $(u_1 - d_1)^{\alpha}(u_2 - d_2)^{1-\alpha}$ and gives
party~1 the share $\alpha$ of the surplus: $u_1 = d_1 + \alpha(V - d_1 - d_2)$.

\begin{keyidea}
Threat points and bargaining weights move the transfers, never the real decision: every
solution sits on the line $u_1 + u_2 = V$, so it uses a value-maximizing $y$. To get a larger
share, improve your outside option ($d_i$) or your bargaining weight ($\alpha$), not the size
of $V$. Section~\ref{sec:b-costs} shows why that private incentive is socially wasteful.
\end{keyidea}

% ---------------------------------------------------------------------------------------
\section{Alternating offers}\label{sec:b-rub}

\subsection{Finite horizon (MWG Example 9.AA.1; Stahl 1972)}

Two players divide $v$ dollars. In odd periods player~1 offers a split and player~2 accepts or
rejects; in even periods the roles swap. A dollar received in period $t$ is worth $\delta^{t-1}$
today, $\delta \in (0,1)$. If no agreement is reached by period $T$, both get nothing
(MWG, p.~296).

\begin{prop}{Proposition 3 (finite horizon)}
The game has a unique subgame perfect Nash equilibrium (SPNE). Agreement is reached at once, and
the first proposer receives
\[
  v_1(T) = v\,\big[1 - \delta + \delta^2 - \dots + (-\delta)^{T-1}\big]
         = v\,\frac{1 - (-\delta)^{T}}{1 + \delta},
\]
which is $v\,(1 + \delta^T)/(1+\delta)$ for odd $T$ and $v\,(1 - \delta^T)/(1+\delta)$ for even $T$.
Player~2 receives $v - v_1(T)$.
\end{prop}

\begin{proof}
Backward induction. In period $T$ the responder accepts any offer (rejecting gives zero), so the
proposer takes everything: share $s_1 = 1$. In period $T-1$ the responder is the next proposer;
rejecting is worth $\delta \cdot s_1$ to him or her in period-$(T-1)$ terms, so the proposer
offers exactly that and keeps $s_2 = 1 - \delta s_1$. In general $s_{k+1} = 1 - \delta s_k$ with
$s_1 = 1$, whose solution is $s_k = \sum_{j=0}^{k-1}(-\delta)^j$. At $k = T$ this is the formula.
Uniqueness: in each subgame the responder accepts iff the offer is at least the continuation
value (accepting at indifference, as MWG and Varian assume), and the proposer's best offer is
pinned down.
\end{proof}

Two lessons. \emph{Who proposes last matters}: with $T = 1$ the proposer takes everything; with
$T = 2$ the second mover gets $\delta v$. \emph{The horizon matters less and less}: as
$T \to \infty$, $v_1(T) \to v/(1+\delta)$ whichever player moves last. With $\delta = 0.6$:
$v_1(3) = 1 - 0.6 + 0.36 = 0.76$, and $v_1(T) \to 0.625$.

\begin{example}{Varian's three-day game (Varian, pp.~581--582)}
Alice and Bob divide \$1 over at most three days; Alice offers on days 1 and 3, Bob on day 2.
Alice discounts at $\alpha$ per day, Bob at $\beta$. On day 3 Alice takes everything. On day 2,
Bob must offer Alice $\alpha$ (the value of \$1 tomorrow) and keeps $1 - \alpha$. On day 1, Alice
must offer Bob $\beta(1 - \alpha)$ and keeps
\[
  1 - \beta(1 - \alpha).
\]
Each party's share rises with its \emph{own} patience: a patient Alice ($\alpha$ near 1) leaves
Bob little to protect on day 2.
\end{example}

\subsection{Infinite horizon: Rubinstein (MWG Example 9.AA.2)}

Remove the deadline and let the players' discount factors differ, $\delta_1, \delta_2 \in (0,1)$.

\begin{prop}{Proposition 4 (Rubinstein 1982)}
The infinite-horizon alternating-offers game has a unique SPNE. Player~1 offers at once
\[
  x_1^* = \frac{1 - \delta_2}{1 - \delta_1\delta_2}, \qquad x_2^* = \frac{\delta_2(1 - \delta_1)}{1 - \delta_1\delta_2},
\]
and player~2 accepts. With $\delta_1 = \delta_2 = \delta$, the shares are $1/(1+\delta)$ and
$\delta/(1+\delta)$, the limit of Proposition~3. In Varian's notation (Alice $\alpha$, Bob
$\beta$): Alice $(1-\beta)/(1-\alpha\beta)$, Bob $\beta(1-\alpha)/(1-\alpha\beta)$ (Varian, p.~582).
\end{prop}

\begin{proof}[Proof sketch (Shaked and Sutton's bounding argument, used by MWG)]
Let $M_1$ and $m_1$ be the supremum and infimum of player~1's SPNE payoffs in any subgame where
he or she proposes, and $M_2, m_2$ the same for player~2 when proposing. The game looks the same
from every proposing node, so these bounds are stationary.
\begin{itemize}
  \item Player~2 rejects anything below $\delta_2 m_2$ and accepts anything above $\delta_2 M_2$.
  So player~1 gets at least $1 - \delta_2 M_2$ and at most $1 - \delta_2 m_2$ (delay only shrinks
  the pie): $m_1 \ge 1 - \delta_2 M_2$ and $M_1 \le 1 - \delta_2 m_2$.
  \item Symmetrically, $m_2 \ge 1 - \delta_1 M_1$ and $M_2 \le 1 - \delta_1 m_1$.
\end{itemize}
Combine: $M_1 \le 1 - \delta_2(1 - \delta_1 M_1)$, so $M_1 \le (1-\delta_2)/(1-\delta_1\delta_2)$;
and $m_1 \ge 1 - \delta_2(1 - \delta_1 m_1)$, so $m_1 \ge (1-\delta_2)/(1-\delta_1\delta_2)$. Hence
$M_1 = m_1 = x_1^*$, and likewise for player~2. The strategies ``always offer the other
$\delta_j x_j^*$, accept any offer of at least $\delta_i x_i^*$'' attain it, and immediate acceptance
follows because any delay loses value both players could share.
\end{proof}

\begin{figure}[ht]
\centering
\begin{tikzpicture}
\begin{axis}[width=0.7\linewidth, height=5.2cm, xlabel={horizon $T$},
  ylabel={first proposer's share}, ymin=0.3, ymax=1.05, xmin=0.5, xmax=12.5,
  xtick={1,...,12}, grid=major, grid style={gray!20}, legend pos=north east,
  legend style={font=\footnotesize}]
\addplot[only marks, mark=*, defc] coordinates {(1,1) (2,0.4) (3,0.76) (4,0.544) (5,0.6736)
  (6,0.59584) (7,0.642496) (8,0.6145024) (9,0.63129856) (10,0.621220864) (11,0.627267482)
  (12,0.623639511)};
\addlegendentry{$v_1(T)$, $\delta = 0.6$}
\addplot[dashed, warnc, domain=0.5:12.5] {0.625};
\addlegendentry{Rubinstein $1/(1+\delta) = 0.625$}
\end{axis}
\end{tikzpicture}
\caption{Finite-horizon shares oscillate (the last proposer gains) and converge to the
infinite-horizon share.}
\end{figure}

\textbf{What the model says about bargaining power.}
\begin{enumerate}
  \item \emph{Agreement is immediate and efficient.} Under complete information no value is lost
  to delay: the threat of delay is what divides the surplus, but it is never carried out.
  \item \emph{Patience is power.} $x_1^*$ rises in $\delta_1$ and falls in $\delta_2$. If
  $\delta_1 \to 1$ with $\delta_2 < 1$, player~1 takes everything (Varian, p.~582).
  \item \emph{The first-mover advantage vanishes as offers speed up.} Write
  $\delta_i = e^{-r_i \Delta}$, with $\Delta$ the time between offers. As $\Delta \to 0$,
  $x_1^* \to r_2/(r_1 + r_2)$: the division depends only on relative impatience. With equal
  rates it tends to $1/2$, the symmetric Nash split, which links the strategic model to the
  axiomatic one of \S\ref{sec:b-nash} (Binmore, Rubinstein and Wolinsky 1986; not in the course
  books).
  \item \emph{The real decision is untouched.} Discount factors move $x$, never $y$: the
  alternating-offers game is a model of the transfers in the Coase theorem.
\end{enumerate}

\textbf{Computing it.} The finite-horizon recursion is three lines of Python; it reproduces the
figure and, with $T$ large, Proposition~4 (\texttt{check\_part2.py} uses it):
\begin{verbatim}
def proposer_share(T, d1, d2):    # pie 1; player 1 proposes first
    share = 1.0                   # the last proposer takes everything
    who = (T - 1) % 2             # last proposer: 0 = player 1, 1 = player 2
    for _ in range(T - 1):        # one period back: the responder (who)
        d = d1 if who == 0 else d2   # must get the discounted value of
        share = 1 - d * share        # proposing next period
        who = 1 - who
    return share
\end{verbatim}

% ---------------------------------------------------------------------------------------
\section{Commitment, reputation and the ultimatum game}\label{sec:b-commit}

In alternating offers, power comes from what a player can credibly do next. \emph{Commitment}
changes that by removing options. \MR's examples \pg{133}: in 1066 William the Conqueror burned
his fleet behind his army, and Cort\'es repeated the stratagem in Mexico; the commitment ``deterred
the opposition, who then themselves retreated''. When Apple built a highly specialized plant for
the Macintosh, it made dropping the product more expensive, and so committed itself in the eyes
of employees, competitors and customers. The cost is real: noncontractual commitment ``may be
worthwhile, but costly'': the fleets were perfectly good, and the specialized plant may become
obsolete sooner \pg{139}.

\textbf{Reputation as a substitute} \pg{139--140}. Concern for reputation checks opportunism when
contracts cannot. Its value grows with the frequency of similar transactions, the horizon over
which they recur, and their profitability. Hence, in relational contracts, discretion over
unforeseen events should go to the party with the most to lose from a damaged reputation: the
one with the longer horizon, more visibility, greater size and more frequent transactions,
which in employment is usually the firm, not the employee \pg{140}. Bargaining with terrorists
or renegotiating executive pay advertises that there is no commitment \pg{139}.

\textbf{The ultimatum game} (Varian, pp.~582--583). Alice proposes a division of \$1; if Bob
accepts it is implemented, otherwise both get nothing. The SPNE gives Alice almost everything
(``something like 99 cents''). In experiments, the average offer by U.S. undergraduates is about
45 cents and is usually accepted. Proposers are close to maximizing expected payoff given the
observed rejection rates; it is the responders who depart from the theory, rejecting small offers
at a cost to themselves. Explanations: social norms of fair division, which differ across
cultures, and a taste for punishing unfair offers.

\begin{keyidea}
A one-shot offer is commitment in its purest form: the proposer is committed to walking away.
A responder who could credibly commit to rejecting anything below $k$ would receive $k$. A
known taste for fairness or vengeance acts as exactly such a commitment, which is the question
\MR{} pose in FFT~5.1 (pp.~163--164); FFT~5.2 (p.~164) asks what makes a signal of intent
credible.
\end{keyidea}

% ---------------------------------------------------------------------------------------
\section{Private information: when efficient trade is infeasible}\label{sec:b-private}

\MR's turn is one sentence: efficiency is defined relative to the feasible outcomes, and
``what is feasible---and thus what is efficient---turns out to depend crucially on the
informational conditions'' \pg{140}. With private information, an agreement may not be reached
even when it would be efficient, and if it is reached it is not typically efficient by
full-information standards, ``especially when the parties have an option not to participate''
\pg{140}.

\subsection{\MR's example: bargaining over a sale}

A seller owns one unit; a buyer wants it. Each knows only his or her own value \pg{140--141}:

\begin{center}
\begin{tabular}{lcc}
\toprule
 & Seller's value \$0 (prob.\ 0.8) & Seller's value \$2 (prob.\ 0.2) \\
\midrule
Buyer's value \$1 (prob.\ 0.2) & Trade & No trade \\
Buyer's value \$3 (prob.\ 0.8) & Trade & Trade \\
\bottomrule
\end{tabular}

\smallskip
{\footnotesize Efficient outcomes (\MR{} Table 5.2, p.~141). Trade is efficient with probability
$1 - 0.2 \times 0.2 = 0.96$.}
\end{center}

\textbf{The temptation.} A buyer worth \$3 would like to claim \$1 to pay less; a seller worth
\$0 would like to claim \$2 to be paid more. Misrepresentation risks the sale falling through,
but ``this may be a risk worth taking'' \pg{141}.

\textbf{The mediator device.} Ask whether a neutral mediator who sets prices as a function of the
two reports could make honest reporting optimal \emph{and} trade happen exactly when efficient.
If no mediator can, the parties cannot do it by themselves either \pg{141}. Two kinds of
constraints bind:
\begin{itemize}
  \item \emph{participation constraints}: nobody can be forced to trade at a loss given the
  value he or she reports;
  \item \emph{incentive compatibility constraints}: each type must do at least as well reporting
  truthfully as reporting the other type.
\end{itemize}
The surplus a type must be left to report honestly is its \emph{informational rent} (n.~7, p.~141).

\textbf{Pinning down the prices} (\MR{} Table 5.3, p.~142). A seller who claims \$2 must get at
least \$2, and paying more only rewards the claim: price \$2. A buyer who claims \$1 must pay at
most \$1: price \$1. One price is left, $p$, for reports (seller \$0, buyer \$3).

\begin{center}
\begin{tabular}{lcc}
\toprule
 & Buyer reports \$1 & Buyer reports \$3 \\
\midrule
Seller reports \$0 & price \$1 & price $p$ \\
Seller reports \$2 & no trade & price \$2 \\
\bottomrule
\end{tabular}
\end{center}

\emph{Seller worth \$0} (the buyer reports honestly):
\[
  \underbrace{0.2 \times 1 + 0.8\,p}_{\text{honest}} \;\ge\;
  \underbrace{0.2 \times 0 + 0.8 \times 2}_{\text{claim \$2}} = 1.60
  \quad\Longrightarrow\quad p \ge 1.75 .
\]
\emph{Buyer worth \$3}:
\[
  \underbrace{0.2 \times (3 - 2) + 0.8\,(3 - p)}_{\text{honest}} = 2.60 - 0.8\,p \;\ge\;
  \underbrace{0.8 \times (3 - 1)}_{\text{claim \$1}} = 1.60
  \quad\Longrightarrow\quad p \le 1.25 .
\]
No $p$ satisfies both. ``It is impossible to find any prices that always make it in the parties'
individual interests to report truthfully and to trade whenever trade is, in fact, value
increasing'' \pg{142}.

\begin{defn}{The revelation principle (\MR{} n.~8, pp.~142--143)}
Any pattern of outcomes achievable by any mechanism that respects incentive constraints can be
achieved by a mechanism in which the parties report their values and are given incentives to
report honestly. Hence restricting attention to honest reporting, as above, loses no generality:
the impossibility holds for every bargaining procedure.
\end{defn}

\textbf{Reverse the probabilities} \pg{143}. Let the buyer be worth \$1 with probability 0.8 and
the seller \$2 with probability 0.8, and set $p = 1.50$. A seller worth \$0 gets
$0.8 \times 1 + 0.2 \times 1.50 = 1.10$ by reporting honestly and $0.2 \times 2 = 0.40$ by
claiming \$2. A buyer worth \$3 gets $0.8 \times 1 + 0.2 \times 1.50 = 1.10$ honestly and
$0.2 \times 2 = 0.40$ by claiming \$1. Both constraints hold (checking the low types and participation
too, any $p$ between 0 and 3 works), so efficient trade is attainable.

\begin{keyidea}
Lying about your value is cheap when trade is likely to happen anyway, and expensive when it puts
a likely trade at risk. In the original example the high-value buyer and the low-value seller
each face a partner who is probably willing to trade (probability 0.8), so exaggerating costs
little. Reverse the odds and every misreport kills a trade 80\% of the time \pg{143}.
\end{keyidea}

The example is a discrete instance of a general result for bilateral trade with independent
private values whose supports overlap: no mechanism is efficient, incentive compatible, individually
rational and budget balanced (Myerson and Satterthwaite 1983; the name is not in the course
books).

% ---------------------------------------------------------------------------------------
\section{Incentive efficiency}\label{sec:b-ie}

If incentive constraints cannot be escaped, failing the full-information benchmark is not a
failure of the parties. \MR{} change the benchmark \pg{143}:

\begin{defn}{Incentive efficiency}
A mechanism is \emph{incentive efficient} if no alternative mechanism that is feasible in the
presence of the incentive constraints is preferred by all the relevant parties. By the efficiency
principle, bargainers should be expected to find incentive-efficient mechanisms.
\end{defn}

\subsection{Split-the-difference bargaining}

The buyer's value $B$ and the seller's value $S$ are independent and uniform on $[0,1]$, each
private. Both name a price at once; if the seller's ask $s$ is below the buyer's bid $b$, they
trade at $(s + b)/2$; otherwise both walk away \pg{143}.

\textbf{Small lies pay} \pg{144}. Let the buyer bid truthfully and a seller with $S = 0.50$ ask
$0.52$. Trades are lost only when $B \in (0.50, 0.52)$, probability $0.02$, each losing at most
$\$0.01$ and $\$0.005$ on average: expected loss $0.02 \times 0.005 = \$0.0001$. Against that, the
price rises by $\$0.01$ whenever $B \ge 0.52$ (probability $0.48$): gain $\$0.0048$. Net gain
$\$0.0047$. The losses from misreporting fall on trades that ``were hardly profitable anyway''.

\begin{prop}{Proposition 5 (the linear equilibrium, \MR{} n.~9, p.~144)}
There is an equilibrium in linear strategies
\[
  s(S) = \tfrac14 + \tfrac23 S, \qquad b(B) = \tfrac1{12} + \tfrac23 B,
\]
and trade occurs if and only if $B > S + \tfrac14$.
\end{prop}

\begin{proof}
Suppose the buyer bids $b = \beta_1 + \beta_2 B$. A seller with value $S$ who asks $s$ trades when
$B \ge (s - \beta_1)/\beta_2$ and earns
\[
  \pi_S(s) = \int_{(s-\beta_1)/\beta_2}^{1} \Big(\frac{s + \beta_1 + \beta_2 B}{2} - S\Big)\, dB .
\]
Differentiating (Leibniz rule; at the lower limit the price equals $s$):
$\pi'_S(s) = \tfrac12\big(1 - \tfrac{s - \beta_1}{\beta_2}\big) - \tfrac{1}{\beta_2}(s - S) = 0$,
so $s = \tfrac23 S + \tfrac{\beta_1 + \beta_2}{3}$. Symmetrically, if the seller asks
$\sigma_1 + \sigma_2 S$, a buyer with value $B$ who bids $b$ trades when $S \le (b - \sigma_1)/\sigma_2$
and the first-order condition gives $b = \tfrac23 B + \tfrac{\sigma_1}{3}$. Matching
coefficients: $\sigma_2 = \beta_2 = \tfrac23$, $\sigma_1 = (\beta_1 + \tfrac23)/3$,
$\beta_1 = \sigma_1/3$, so $\sigma_1 = \tfrac14$ and $\beta_1 = \tfrac1{12}$. Finally
$b > s \iff \tfrac1{12} + \tfrac23 B > \tfrac14 + \tfrac23 S \iff B > S + \tfrac14$.
\end{proof}

\textbf{How much is lost.} Let $D = B - S$; on $[0,1]$ its density is $1 - D$. Efficient trade
happens with probability $\tfrac12$ and creates expected gains $\int_0^1 D(1-D)\,dD = \tfrac16$.
In equilibrium trade happens with probability $\int_{1/4}^{1}(1 - D)\,dD = \tfrac{9}{32}$, and the
realized gains are $\int_{1/4}^{1} D(1-D)\,dD = \tfrac{9}{64}$, which is $\tfrac{27}{32} \approx 84\%$
of the efficient total. Trade is cut almost in half, but surplus by only 16\%, because the trades
lost are the ones with small gains.

\begin{figure}[ht]
\centering
\begin{tikzpicture}[scale=4.2]
  \fill[exc!18] (0,0.25) -- (0,1) -- (0.75,1) -- cycle;
  \fill[warnc!18] (0,0) -- (0,0.25) -- (0.75,1) -- (1,1) -- cycle;
  \draw (0,0) rectangle (1,1);
  \draw[thick, gray] (0,0) -- (1,1);
  \draw[thick, exc!70!black] (0,0.25) -- (0.75,1);
  \node[font=\footnotesize, align=center] at (0.2,0.8) {trade\\$B > S + \frac14$};
  \node[font=\footnotesize, rotate=45] at (0.55,0.675) {lost};
  \node[font=\footnotesize, align=center] at (0.7,0.25) {no gains\\$B < S$};
  \node[below] at (0.5,0) {\footnotesize seller's value $S$};
  \node[rotate=90, above] at (0,0.5) {\footnotesize buyer's value $B$};
  \node[below, font=\scriptsize] at (0,0) {0}; \node[below, font=\scriptsize] at (1,0) {1};
  \node[left, font=\scriptsize] at (0,1) {1};
\end{tikzpicture}
\caption{Split-the-difference with uniform values. Efficient trade: everything above the
diagonal (probability $1/2$). Equilibrium trade: the green triangle (probability $9/32$). The red
band $S < B \le S + \frac14$ holds the efficient trades lost to the incentive constraints.}
\end{figure}

\MR{} add that these rules ``actually maximize efficiency subject to the constraints imposed by
the private information'' \pg{144}: no other procedure does better for both. Letting the
parties continue to bargain does not help; the inefficiency then shows up ``either in the form
of costly disagreements or as a costly delay until an agreement is reached'' \pg{144}.
The general pattern: when incentive constraints matter, ``trade takes place only if the gains
from trade are sufficiently large'' \pg{144}.

% ---------------------------------------------------------------------------------------
\section{Many participants: the public project}\label{sec:b-many}

The problem worsens as more people must agree \pg{145}. A project costs \$1 in total; each of $N$
people values it at \$2 with probability $p$ or at \$0, privately and independently. Anyone can
drop out (and lose the benefit), so nobody pays more than \$2 and someone claiming \$0 pays
nothing. Value maximization says: build whenever anyone values it. If $m \ge 1$ people report
\$2, build and split the cost among them.

\textbf{$N = 2$.} A person worth \$2 compares honest reporting with claiming \$0, given the other
reports honestly \pg{145}:
\[
  \underbrace{p\,(2 - 0.5) + (1 - p)(2 - 1)}_{\text{honest}} \;\ge\; \underbrace{p \cdot 2 + (1-p)\cdot 0}_{\text{free ride}}
  \quad\Longleftrightarrow\quad 1 + 0.5\,p \ge 2p \quad\Longleftrightarrow\quad p \le \tfrac23 .
\]
If $p > 2/3$ the other is too likely to pay alone, and the temptation to free ride wins
\pg{146}. The remedy is to build less often when the reported gains look small: build with
probability $q$ when only one person reports \$2. The constraint becomes
$1.5\,p + (1-p)\,q \ge 2pq$, i.e.\ $q \le \dfrac{3p}{2(3p - 1)}$, which falls to $\tfrac34$ as
$p \to 1$ \pg{146}: a quarter of the time a volunteer appears, the project is not built.

\textbf{$N = 3$.} With honest others, the constraint
\[
  p^2\big(2 - \tfrac13\big) + 2p(1-p)\big(2 - \tfrac12\big) + (1-p)^2(2 - 1) \;\ge\; 2\,\big[1 - (1-p)^2\big]
\]
holds only for $p \le 0.4417$ (\MR: ``less than about 0.44'', p.~146). At $p = \tfrac12$, building
with probability $q$ when exactly one person reports \$2 requires
$\tfrac14 \cdot \tfrac53 + \tfrac12 \cdot \tfrac32 + \tfrac14 q \ge \tfrac14 \cdot 2 + \tfrac12 \cdot 2q$,
i.e.\ $q \le \tfrac89$: the project goes ahead ``less than 90 percent of the time'' when one
pivotal person is willing to pay for all of it \pg{146}.

\textbf{Large $N$.} Lying is punished only when you are \emph{pivotal}: when your report decides
whether the project is built. As $N$ grows it becomes nearly certain that someone else values the
project, so the cost of free riding vanishes and misrepresentation cannot be prevented
\pg{146}. Allowing for everyone's strategic reporting, \MR{} report that the probability that
someone volunteers tends to ``approximately two chances in seven, even though it is essentially
certain that funding the project is value maximizing'' (pp.~146--147).

\begin{keyidea}
More participants make agreement harder, not easier. The free-rider problem is a private
information problem: if values were known, the group could assign costs and build whenever
efficient.
\end{keyidea}

% ---------------------------------------------------------------------------------------
\section{Bargaining costs and investments in bargaining position}\label{sec:b-costs}

\textbf{Bargaining costs} are ``the transaction costs of reaching mutually acceptable
agreements'' \pg{147}. Bounded rationality imposes them even without strategic behaviour: it takes
time to list contingencies, find efficient plans and settle divisions, and the resulting
incomplete contracts create further costs of achieving commitment. Private information adds
breakdown and delay. Hence ``simple value maximization computed without regard to the
constraints imposed by incentives is a misleading basis for prediction'' \pg{147}.

\textbf{Measurement costs} \pg{147--148}. When the quality of a good is unknown to both sides
(oil under a tract of land, the durability of equipment), each side can gain a bargaining
advantage by measuring it. If one side measures, the other must too, or face adverse selection.
Yet the measurement changes only the price, not who ends up with the good: if each side spends
$k$, the good still changes hands and total wealth falls by $2k$. ``The costs of gathering that
information are subtracted from the parties' total wealth; there is no offsetting benefit unless
the information affects a productive decision or the allocation of the asset'' \pg{147}.

\begin{example}{De Beers (\MR{} p.~148)}
De Beers, which controlled over 80\% of world diamond supply, sold rough stones in ``sights'':
packets sorted on gross characteristics only, offered at a formula price on a strict
take-it-or-leave-it basis, with no negotiation over price or contents, and with exclusion from
future sales for any buyer who refused a sight. A monopoly reading explains part of this. The
value-maximizing reading explains the details. If bargaining were allowed, buyers would examine
each stone closely to argue over price, and De Beers would have to do the same in self-defence;
only the final buyer, who cuts the stone, needs that examination. Forbidding negotiation removes
the incentive to measure; exclusion stops ``leave it'' from being used as a bargaining threat;
and buyers know the stones have not been picked over. Both sides save. The same logic is offered
for block-booking of films, prepackaged produce, product warranties and percentage royalties
for authors \pg{148}.
\end{example}

\textbf{Investing in bargaining advantages} \pg{149}. Measurement costs belong to a larger
class: expenditures with private value because they shift the division, but no social value.
\MR's example is weapons systems and the Strategic Defense Initiative funded as ``bargaining
chips'' to be traded away in arms-control talks. Haggling, posturing and delay are wastes of the
same kind, and if the parties foresee them they may never attempt the agreement at all.

\begin{example}{Preview of Topic 6: a specific investment and an ex-post split}
A supplier can make a relationship-specific investment $e \ge 0$ at private cost $e^2/2$ that
raises the joint value by $a\,e$. Contracts cannot specify $e$, and the value is split 50/50
by bargaining after the investment (the symmetric Nash solution of \S\ref{sec:b-nash}, with the
investment sunk). The supplier maximizes $\tfrac12 a e - \tfrac12 e^2$ and picks $e = a/2$, while
the efficient level maximizes $a e - \tfrac12 e^2$ and is $e = a$. Total surplus is
$\tfrac38 a^2$ instead of $\tfrac12 a^2$: a quarter is lost. Ex-post bargaining is efficient given
$e$ (Coase) but distorts the investment that came before it. This is the hold-up problem
(\MR{} Ch.~5, investments and specific assets, pp.~134--139; Varian ``Hold Up'', p.~579), which
returns in Aula~9.
\end{example}

% ---------------------------------------------------------------------------------------
\section{How the pieces fit}

\begin{center}
\begin{tabularx}{\linewidth}{>{\raggedright\arraybackslash}p{3.4cm}X>{\raggedright\arraybackslash}p{3.6cm}}
\toprule
Hypothesis & What happens when it fails & Where \\
\midrule
No wealth effects & the real choice depends on who holds the rights (WTP $\neq$ WTA) & \S\ref{sec:b-wealth}, \pg{39} \\
Unlimited transfers & value-creating deals blocked by limited capital & \pg{39} \\
Complete information & efficient trade may be infeasible; incentive-efficient rules drop small-gain trades & \S\ref{sec:b-private}--\S\ref{sec:b-ie} \\
Few parties & free riding grows with $N$; volunteering $\to$ about $2/7$ & \S\ref{sec:b-many} \\
Costless bargaining & measurement, haggling, bargaining chips, delay & \S\ref{sec:b-costs} \\
Contractible actions & ex-post bargaining distorts investment (hold-up) & \S\ref{sec:b-costs} \\
\bottomrule
\end{tabularx}
\end{center}

The Coase theorem is a benchmark: under its hypotheses organization would not matter, because
any arrangement would be renegotiated to the efficient one. Each failure in the table is a
transaction cost, and the rest of the course is about the organizations designed to reduce
them: the coordination and motivation problems of Part~I, seen from the contract side.

% ---------------------------------------------------------------------------------------
\section{Common mistakes}

\begin{center}
\begin{tabularx}{\linewidth}{>{\raggedright\arraybackslash}p{5.2cm}X}
\toprule
Mistake & Correction \\
\midrule
``Bargaining power decides what gets done'' & Under no wealth effects it decides only the transfers $x$; efficiency decides $y$ (pp.~36--38). \\
``Coase: property rights do not matter'' & They decide the distribution, and with wealth effects the allocation too (WTP/WTA example). \\
``Value maximization = efficiency, always'' & Only without wealth effects; otherwise sum-maximizing is sufficient, not necessary (n.~10, p.~37). \\
``Rational parties always reach efficient deals'' & With private information the efficient outcome can violate incentive constraints (pp.~141--142). \\
``A mechanism that misses some gains is inefficient'' & Judge it against feasible mechanisms: incentive efficiency (p.~143). \\
``The revelation principle says people tell the truth'' & It says honest-reporting mechanisms lose no generality (n.~8). \\
``The first mover always wins'' & In Rubinstein the advantage vanishes as offers speed up; patience decides (\S\ref{sec:b-rub}). \\
``More participants make agreement easier'' & Pivotality, and with it the cost of free riding, vanishes as $N$ grows (pp.~145--147). \\
``Information gathering is productive'' & Not when it only moves the price: measurement costs are waste (p.~147). \\
``De Beers' rules are just monopoly power'' & They also cut wasteful measurement on both sides (p.~148). \\
\bottomrule
\end{tabularx}
\end{center}

% ---------------------------------------------------------------------------------------
\section{Exercises (number and page only)}

From \texttt{Listas/pratica-aula-01-02.md}; statements are in the books. Suggested order: top to
bottom.

\begin{center}
\begin{tabularx}{\linewidth}{llrX}
\toprule
Exercise & Book & Page & Practises \\
\midrule
32 RQ 3 & Varian & 628 & contract curve vs.\ the outcome of bargaining \\
35 RQ 1--2 & Varian & 687 & property rights, Coase with quasilinear preferences \\
5.1 Front yards & Cabral & 117 & a two-sided externality and bargaining \\
Q 2.1--2.2 & \MR & 53 & value maximization; cost sharing and side payments \\
Q 2.3--2.4 & \MR & 53 & proof of the principle; no wealth effects $\Rightarrow$ quasilinear \\
7.6 Hernan Cortez & Cabral & 180 & commitment \\
7.14 Centipede & Cabral & 182 & backward induction \\
FFT 5.1--5.2 & \MR & 163--164 & fairness preferences; credible commitment \\
9.B.13, 9.B.12 & MWG & 302--303 & finite horizon; Rubinstein with $\delta_1 \ne \delta_2$ \\
Q 5.4--5.5 & \MR & 165 & bargaining with private values, honest-reporting prices \\
FFT 2.5 & \MR & 52 & asset specificity and bargaining position \\
13.10 & Cabral & 346 & ex-post bargaining and investment (hold-up) \\
\bottomrule
\end{tabularx}
\end{center}

% ---------------------------------------------------------------------------------------
\section*{Sources for Part II}
\addcontentsline{toc}{section}{Sources for Part II}

\begin{itemize}
  \item Milgrom, P. \& Roberts, J. (1992). \emph{Economics, Organization, and Management}.
  Prentice-Hall. Ch.~2 (pp.~35--39), Ch.~5 (pp.~133--149). Read from the converted text on
  2026-10-13.
  \item Varian, H. (2014). \emph{Intermediate Microeconomics with Calculus}, 1st ed. Norton.
  \S30.7 and ``The Ultimatum Game'' (pp.~581--583), \S35.2 (p.~670). The syllabus cites the
  8th edition of the non-calculus book; pages may differ.
  \item Mas-Colell, A., Whinston, M. \& Green, J. (1995). \emph{Microeconomic Theory}. OUP.
  Ch.~9, Appendix~A (pp.~296--299).
  \item Cabral, L. (2017). \emph{Introduction to Industrial Organization}, 2nd ed. MIT Press.
  \S5.1 (pp.~96--99).
  \item Outside the course books, cited for context only: Nash (1950), \emph{Econometrica}
  18(2), 155--162; Rubinstein (1982), \emph{Econometrica} 50(1), 97--109; Binmore, Rubinstein \&
  Wolinsky (1986), \emph{RAND Journal of Economics} 17(2), 176--188; Myerson \& Satterthwaite
  (1983), \emph{Journal of Economic Theory} 29(2), 265--281.
  \item Every number in this Part is checked by \texttt{check\_part2.py} (sympy and numpy).
\end{itemize}
